\begin{figure*}[h!]
    \centering
    \includegraphics[width=\linewidth]{figs/assets/capabilities.pdf}
    \caption{Qualitative examples of learned capabilities. We present generated videos demonstrating the model's ability to simulate complex game dynamics. The rows illustrate: (1) fine-grained state tracking, specifically inventory counter synchronization following block placement; (2) global environmental consistency, shown by the simultaneous onset of rain; (3) inventory active item synchronization, torch placement, and generating accurate mining animations; and (4) coherent PVP on complex terrain.}
    \label{fig:capabilities}
\end{figure*}

\section{Experiments}
\label{sec:experiments}
\begin{table}[t]
\centering

\resizebox{\textwidth}{!}{%
\begin{tabular}{lcccccccccc}
\toprule
 & \multicolumn{2}{c}{\textbf{Movement}} 
 & \multicolumn{2}{c}{\textbf{Grounding}} 
 & \multicolumn{2}{c}{\textbf{Memory}} 
 & \multicolumn{2}{c}{\textbf{Building}} 
 & \multicolumn{2}{c}{\textbf{Consistency}} \\
\cmidrule(lr){2-3} 
\cmidrule(lr){4-5} 
\cmidrule(lr){6-7} 
\cmidrule(lr){8-9} 
\cmidrule(lr){10-11}
\textbf{Method} 
& \textbf{VLM} $\uparrow$ & \textbf{FID} $\downarrow$ 
& \textbf{VLM} $\uparrow$ & \textbf{FID} $\downarrow$ 
& \textbf{VLM} $\uparrow$ & \textbf{FID} $\downarrow$ 
& \textbf{VLM} $\uparrow$ & \textbf{FID} $\downarrow$ 
& \textbf{VLM} $\uparrow$ & \textbf{FID} $\downarrow$ \\
\midrule
Frame concat & \textbf{77.1 $\pm$ 2.9} & 68.9 & 53.1 $\pm$ 0.0 & 66.6 & \textbf{37.5 $\pm$ 2.6} & 74.4 & 0.0 $\pm$ 0.0 & 103.2 & 49.5 $\pm$ 5.3 & 129.4 \\
\solaris w/o pretrain  & 69.3 $\pm$ 0.7 & 42.5 & 29.2 $\pm$ 1.5 & 49.9 & 18.8 $\pm$ 0.0 & 67.8 & 0.0 $\pm$ 0.0 & 86.6 & 49.5 $\pm$ 3.2 & 121.4 \\
\solaris & 68.2 $\pm$ 2.7 & \textbf{38.5} & \textbf{62.5 $\pm$ 0.0} & \textbf{38.0} & \textbf{37.5 $\pm$ 2.6} & \textbf{55.1} & \textbf{20.8 $\pm$ 1.5} & \textbf{83.6} & \textbf{71.4 $\pm$ 7.1} & \textbf{99.4} \\
\bottomrule
\end{tabular}%
} %

\caption{Quantitative comparison across tasks. We compare our method against concatenating player observations along the channel dimension following Enigma Multiverse~\cite{enigma2025multiverse} and training from scratch without single-player pretraining. ``VLM'' shows the VLM accuracy $\pm$ 1 (estimated) standard deviation. The standard deviation is estimated by repeating the VLM eval 3 times.}
\label{tab:arch_ablation_results}
\end{table}


\subsection{Qualitative Results}
Here we discuss qualitative results from our flagship, \solaris, model. As shown in~\cref{fig:teaser}, our model is capable of simulating complex aspects of Minecraft gameplay, including building, mining, fighting, and multiplayer viewpoint modeling. The action sequences for each of these scenarios come from bot gameplay not present in the training dataset. We also recorded a third-person perspective of the bots, shown in the left and right columns. As can be seen in the end frames from our model, its generations align with the ground truth game state. Examples of advanced model capabilities---including inventory tracking, simulating weather, placing torches, generating animations, and simulating PvP---are presented in~\cref{fig:capabilities}. ~\cref{fig:qualitative} presents qualitative results from our architecture comparison discussed in~\cref{subsec:architecture_experiments}. While baseline methods degrade over time, our model maintains visual fidelity throughout a long-horizon.


\subsection{Architecture Experiments}
\label{subsec:architecture_experiments}


We compare our architecture implementation to the frame concatenation method of Multiverse ~\cite{enigma2025multiverse}, the only existing multiplayer world model prior to this work. We also test the necessity of single-player pretraining by comparing the performance of our method to the variant trained without the single-player model initialization. Our method produces superior visual results both qualitatively, as shown in~\cref{fig:qualitative}, and quantitatively across all evaluation categories, as shown in~\cref{tab:arch_ablation_results}. All architecture variants are strong at action following in motion-based trajectories, and get a high VLM score on the evaluation in the corresponding category, as shown in~\cref{tab:arch_ablation_results}. Our method shows superior performance on difficult scenarios involving building, scene consistency, and player grounding, reflected by the higher VLM scores in those categories. Although the frame concatenation method outperforms in our Movement evaluation, we find qualitatively that it exhibits action hallucinations in the presence of no-op actions. 

\subsection{Self-Forcing Ablations}
\label{subsec:sf_ablations}

We ablate each component of our Self Forcing pipeline. We study the two main stages of CausVid, ODE regression initialization and few-step distillation, and find that straightforward causal finetuning is sufficient instead. Although the original Self Forcing paper assumes the generator to be a few-step model at the start of training, we find that the few-step ability can be learned simultaneously with stable autoregressive generation in Self Forcing. This allows use simple finetuning instead of CausVid.

Next, we study the ability to backpropagate to the KV representations of the self attention layer, which is feasible with the memory savings of our method. Allowing KV backpropagation achieves better visual results than all other variants based on FID, as shown in~\cref{tab:sf_ablation_results}. We do observe that this causes decreased performance in action following for some categories. However, our method remains competitive across all categories and excels in the challenging Building and Consistency VLM tasks.



















\begin{table*}[h!]
\centering

\resizebox{\textwidth}{!}{%
\begin{tabular}{lcccccccccccc}
\toprule
\multicolumn{3}{c}{\textbf{Components}} & 
\multicolumn{2}{c}{\textbf{Movement}} & 
\multicolumn{2}{c}{\textbf{Grounding}} & 
\multicolumn{2}{c}{\textbf{Memory}} & 
\multicolumn{2}{c}{\textbf{Building}} & 
\multicolumn{2}{c}{\textbf{Consistency}} \\
\cmidrule(r){1-3}
\cmidrule(lr){4-5}
\cmidrule(lr){6-7}
\cmidrule(lr){8-9}
\cmidrule(lr){10-11}
\cmidrule(l){12-13}
\textbf{Init.} & \textbf{Pre-DMD} & \textbf{KV-BP} &
\textbf{VLM} $\uparrow$ & \textbf{FID} $\downarrow$ &
\textbf{VLM} $\uparrow$ & \textbf{FID} $\downarrow$ &
\textbf{VLM} $\uparrow$ & \textbf{FID} $\downarrow$ &
\textbf{VLM} $\uparrow$ & \textbf{FID} $\downarrow$ &
\textbf{VLM} $\uparrow$ & \textbf{FID} $\downarrow$ \\
\midrule
ODE Reg & \texttimes & \checkmark & 23.4 $\pm$ 1.5 & 65.3 & 3.1 $\pm$ 0.0 & 56.6 & 0.0 $\pm$ 0.0 & 99.5 & 3.1 $\pm$ 0.0 & 95.7 & 49.0 $\pm$ 4.2 & 142.3 \\
Causal FT & \checkmark & \checkmark & 21.4 $\pm$ 2.7 & 49.1 & 3.1 $\pm$ 0.0 & 40.4 & 0.0 $\pm$ 0.0 & 55.7 & 8.3 $\pm$ 1.5 & 90.5 & 55.2 $\pm$ 3.7 & 160.1 \\
Causal FT & \texttimes & \texttimes & \textbf{78.6 $\pm$ 0.7} & 60.3 & \textbf{72.9 $\pm$ 1.5} & 55.2 & \textbf{49.0 $\pm$ 2.9} & 63.8 & 15.6 $\pm$ 0.0 & 87.4 & 70.8 $\pm$ 6.6 & 105.1 \\
\rowcolor{gray!15} Causal FT & \texttimes & \checkmark & 68.2 $\pm$ 2.7 & \textbf{38.5} & 62.5 $\pm$ 0.0 & \textbf{38.0} & 37.5 $\pm$ 2.6 & \textbf{55.1} & \textbf{20.8 $\pm$ 1.5} & \textbf{83.6} & \textbf{71.4 $\pm$ 7.1} & \textbf{99.4} \\
\bottomrule
\end{tabular}%
}

\caption{Ablations on Self Forcing training variants. We study the initialization of the causal model, finding simple causal finetuning with Diffusion Forcing~\cite{chen2024diffusionforcingnexttokenprediction} to suffice. We also find that doing few-step distillation before Self-Forcing is ineffective. Finally, we find that enabling KV cache backpropagation improves visual quality. ``VLM'' shows the VLM accuracy $\pm$ 1 (estimated) standard deviation. The standard deviation is estimated by repeating the VLM eval 3 times.}
\label{tab:sf_ablation_results}
\end{table*}



